\scanpage{44}\mspage{$\underline{\overline{\mathrm{XVI}}}$}
\uline{Lemme} $\underline{\overline{\mathrm{XVII}}}$. Si $x\in\Gamma$ et $D_x$ est un disque de centre $x$ qui ne coupe que les arêtes auxquelles $x$ appartient alors $D_x\cap U\ne\varnothing$, $D_x\cap V\ne\varnothing$.

\uline{Démonstration} : D'aprés la remarque qui suit la def. 1 il suffit d'exhiber un disque $D_x$ qui ne coupe que les arètes auxquelles $x$ appartient et tel que $D_x\cap U\ne\varnothing$, $D_x\cap V\ne\varnothing$

Si $x\in\Gamma^{1}_{ss'}-\{s,s'\}$ soit $D_x$ un disque de centre $x$ vérifiant : a) $D_x\subset\mathbb{R}^{2}-\Lambda^{2}$ ; b) $D_x$ satisfait b) de la \unsure{definition|defnition} 1 par rapport à $\Lambda^{1}$, donc $D_x\cap\mathrm{ext}\,\Lambda^{1}\ne\varnothing$ et par conséquent $D_x\cap V\ne\varnothing$. Aussi $D_x\subset\mathrm{ext}\,\Lambda^{2}$ et $D_x\cap\mathrm{int}\,\Lambda^{1}\ne\varnothing$ entraîne $D_x\cap V\ne\varnothing$. Si $x\in\Gamma^{2}_{ss'}-\{s,s'\}$ un argument analogue montre que $D_x\cap U\ne\varnothing$ et $D_x\cap V\ne\varnothing$.

Soit maintenant un disque $D_s$ qui satisfait b) de la \unsure{definition|defnition} 1 par rapport a $\Lambda^{1}$ et à $\Lambda^{2}$ on sait par le lemme \inserted{\textcircled{\scriptsize$\underline{\overline{\mathrm{XIV}}}$}} que
\[
\begin{aligned}
D_s\cap{}&\mathrm{ext}\,\Lambda^{1}\ne\varnothing
  &\quad\text{et}\quad D_s\cap{}&(\mathrm{int}\,\Lambda^{1}-\mathrm{int}\,\Lambda^{2})\\[-.5ex]
 &\hphantom{\mathrm{e}}\inserted{\begin{gathered}[t]\cap\\V\end{gathered}}
  &\begin{gathered}[t]\vphantom{\cap}\\\ne\varnothing\end{gathered}
  &\hphantom{(\mathrm{int}\,\Lambda^{1}-}\inserted{\begin{gathered}[t]\cap\\U\end{gathered}}
\end{aligned}
\]
ce qui suffit à prouver le lemme s$\dot{\mathrm{p}}$our $s$, un argument analogue démontre le lemme pour $s'$.

La \unsure{conjonetion|conjonction} du lemme $\underline{\overline{\mathrm{XVI}}}$ et du lemme $\underline{\overline{\mathrm{XVII}}}$ montre le thēorēme 1.

4- \uline{Preuve du théorēme 2} (On suppose que tout \inserted{polygone} $\Gamma$ est un $\unsure{J|j|y}$-polygone)

On procède par récurrence sur $\deleted{\unsure{S|s}}\mathop{\mathrm{card}} S(\Gamma)-2$.

On choisit $s,s'\in S(\Gamma)$ satisfaisant la condition du lemme XI, i.e. $s'$ est visible de $s$ de l'interieur.

\deleted{\unsure{Posons}} \inserted{Soient} $\Gamma^{1}_{ss'}$ et $\Gamma^{2}_{ss'}$ \deleted{\unsure{égale|égals}} les \unsure{deux|dex} arcs de polygone de $\Gamma$ réliant $s$ à $s'$ décrits dans la Remarque 1, et soit $\Lambda_{ss'}$ l'arc de polygone $\Lambda_{ss'}=[s,s']$. Appelons $\Lambda^{1}$ \deleted{l'arc} (resp. $\Lambda^{2}$) le polygone \deleted{f}obtenu à partir de $\Gamma^{1}_{ss'}$ (resp. $\Gamma^{2}_{ss'}$) et $\Lambda_{ss'}$ par \unsure{rēcollement|récollement} selon $\{s,s'\}$.

Il est clair que $\inserted{\mathrm{card}}\,S(\Lambda^{1})<\mathop{\mathrm{card}} S(\Gamma)$ (resp. $\mathop{\mathrm{card}} S(\Lambda^{2})<\mathop{\mathrm{card}} S(\Gamma)$), autre\textsuperscript{.}\\ment
% archive sheet number: 43
